% =============================================================================
% Configuration generale de la campagne "Alan pas riche"
% =============================================================================
\usepackage[T1]{fontenc}
\usepackage[utf8]{inputenc}
\usepackage[french]{babel}
\usepackage{lmodern}
\usepackage{microtype}

\usepackage[a4paper,margin=26mm,headheight=15pt]{geometry}
\usepackage{setspace}
\setstretch{1.06}

\newlength{\recitparskip}
\setlength{\recitparskip}{0.45em plus 0.15em}

\setlength{\parindent}{1.25em}
\setlength{\parskip}{\recitparskip}
\frenchspacing

\widowpenalty=10000
\clubpenalty=10000
\displaywidowpenalty=10000

% --- Titres ------------------------------------------------------------------
\usepackage{titlesec}

\titleformat{\chapter}[display]
  {\normalfont\huge\bfseries}
  {\filleft\Large\chaptername\ \thechapter}
  {1ex}
  {\titlerule\vspace{1.2ex}\filright}

\titlespacing*{\chapter}{0pt}{-15pt}{28pt}

\titleformat{\section}
  {\normalfont\Large\bfseries}
  {\thesection}
  {0.8em}
  {}

% --- En-tetes et pieds de page ------------------------------------------------
\usepackage{fancyhdr}

\pagestyle{fancy}
\fancyhf{}
\fancyhead[L]{\small\nouppercase{\leftmark}}
\fancyhead[R]{\small Alan pas riche}
\fancyfoot[C]{\thepage}

\renewcommand{\headrulewidth}{0.4pt}
\renewcommand{\footrulewidth}{0pt}

\fancypagestyle{plain}{%
  \fancyhf{}
  \fancyfoot[C]{\thepage}
  \renewcommand{\headrulewidth}{0pt}
}

% --- Informations du livre ----------------------------------------------------
\newcommand{\BookTitle}{Alan pas riche}
\newcommand{\BookSubtitle}{Les aventures de Kael, Morrÿs et Obélik}
\newcommand{\BookMeta}{Donjons \& Dragons --- Récit de campagne}
\newcommand{\BookAuthor}{Jérémy Cheynet}

% --- Metadonnees PDF ----------------------------------------------------------
\usepackage[hidelinks]{hyperref}

\hypersetup{
  pdftitle={Alan pas riche - Les aventures de Kael, Morrys et Obelik},
  pdfauthor={\BookAuthor},
  pdfsubject={Recit de campagne - Donjons et Dragons}
}

% --- Tableaux ----------------------------------------------------------------
\usepackage{booktabs}
\usepackage{longtable}
\usepackage{array}
\usepackage{tabularx}

% --- Couleurs et illustrations ------------------------------------------------
\usepackage{xcolor}
\usepackage{graphicx}

\graphicspath{{img/}}

\renewcommand{\figurename}{Figure}
\renewcommand{\listfigurename}{Table des illustrations}

\newcommand{\illustration}[4][0.85\textwidth]{%
  \begin{figure}[htbp]
    \centering
    \includegraphics[width=#1]{#2}
    \caption{#3}
    \label{fig:#4}
  \end{figure}
}

\definecolor{CampaignGray}{gray}{0.25}

\newcolumntype{L}[1]{%
  >{\raggedright\arraybackslash}p{#1}
}

% --- Page de titre ------------------------------------------------------------
\newcommand{\makecampaigntitle}{%
  \begin{titlepage}
    \centering

    \vspace*{0.20\textheight}

    {\Huge\bfseries \BookTitle\par}

    \vspace{1.2em}

    {\Large \BookSubtitle\par}

    \vspace{2.2em}

    {\large\color{CampaignGray} \BookMeta\par}

    \vfill

    {\large \BookAuthor\par}

    \vspace{0.8em}

    {\small Une histoire écrite au fil des parties\par}

    \vspace*{0.08\textheight}
  \end{titlepage}
}

% --- Separation de scenes -----------------------------------------------------
\newcommand{\scenebreak}{%
  \par
  \vspace{1.2em}

  \begin{center}
    \textasteriskcentered
    \quad
    \textasteriskcentered
    \quad
    \textasteriskcentered
  \end{center}

  \vspace{1.2em}
  \par
}

% --- Informations de session -------------------------------------------------
%
% En-tete facultatif au debut d'un chapitre pour conserver la trace
% de la partie dont est issu le recit.
%
% Exemple :
% \sessioninfo{Session d'initiation}{Avant les vacances 2026}
%
\newcommand{\sessioninfo}[2]{%
  \begin{center}
    \small\itshape
    #1%
    \if\relax\detokenize{#2}\relax
    \else
      {} --- #2%
    \fi
  \end{center}

  \vspace{1em}
}

% --- Apartés du narrateur -----------------------------------------------------
\usepackage[most]{tcolorbox}

\newtcolorbox{narrateur}{
    enhanced,
    breakable,
    colback=black!3,
    colframe=black!50,
    boxrule=0pt,
    leftrule=2pt,
    arc=0pt,
    left=10pt,
    right=10pt,
    top=6pt,
    bottom=6pt,
    before skip=10pt,
    after skip=10pt,
    fontupper=\itshape,
    % tcolorbox n'hérite pas du \parskip du document
    before upper={\setlength{\parskip}{\recitparskip}}
}